\section{Endgame}

When one begins to experience the full effect of the ``Times'' we live in, it becomes obvious that there are a multitude of factors at work. Cologero has pointed out the increasing complexity and speed of events in the Kali Yuga. Although at first and superficial glance this would strike one as \emph{entirely} negative (as the rate of ``Change'' at first is so excited as the preclude the inhibitions of right reason and the evidence of unintended and logical consequences) it is also apparent that a particularly high rate of change begins to have the opposite effect -- people quickly, within their lifetimes, see the necessary outcome of events which they misjudged in their youth. The brain no longer has to anticipate the event -- events come full circle within years or decades. This is increasingly the case in our own times -- a good many converts to Tradition come out of the fog of chaos.

It is by no means certain, however, that the full cycle is possible for many. Some simply die along the way, either physically or morally or spiritually. What is contingently different is that in one physical lifetime, many have the opportunity to learn the ins and outs of Chaos simply by ``living'' and surviving long enough to be bludgeoned into Tradition by Reality itself.

The first step in Anti-Tradition was to cast aspersion and doubt upon first principles so that one said with Pilate ``What is Truth?''. In the West, this was both the cynicism of the Renaissance and then the open infidelity of the Enlightenment. Then came the open confusion of Good with Evil, so that each could be transposed: ``Woe unto those who call good evil, and evil, good,'' said Isaiah. In our case, Romanticism signaled the initial onslaught, which culminated in moral relativism in our own day. So what comes after this? Naturally, the third and final stage (which is analogous to Divinization in the Christian tradition, after Purification and Enlightenment) is the utter inability to tell up from down entirely, so that Reality itself is confused with Delusion and Insanity. Because Time has sped up, the moral relativism of the 1900-1960s is quickly passing into the delusional insanity of the post-Clinton years, in which the Psyche becomes impotent to discern what is mere wish-fulfillment and what is actually real. See this blog (for instance\footnote{\url{http://onecosmos.blogspot.com/}}) for an in-depth analysis of this mindset, which came to full bloom in the 1990s. There is a fourth stage after the third, for the simple reason that after the third, when all reserves have been used up in subsidizing the degeneracy of the avant-garde elements, complete and total collapse is the inevitable outcome of these multiple co-morbidities. This is what is spoken of as Chaos.

For an analysis of this process on the purely physical plane, see the Archdruid Report\footnote{\url{http://thearchdruidreport.blogspot.com/}}. But it operates on all planes.

How long this Dark Age (which will be a true Dark Age that makes a mockery out of the propaganda which portrayed the Middle Ages as ``dark'') lasts is up to, quite simply speaking, a spiritual elite.

It is obvious from the increasingly compressed time periods of each stage that we will quickly and shortly (within our lifetimes) enter the final phase. Thus, it is not a mere theoretical imperative to begin to work on the Self, and to plant the seeds of the Future, but a necessity of the highest moral and spiritual order.

Where will we find this order? Will we find it by cobbling together scraps of Nietzsche with a new populism? Here is what the Russian Church\footnote{\url{http://english.pravda.ru/news/russia/29-10-2013/126027-russian\_orthodox\_church-0/}} says: we should not succumb to ``a fruitless populist radicalism'', in order to fight the rising threats of immigration and multi-culturalism. Multiculturalism itself exists globally outside the ivory towers of Western decay only as local cultures flattened by globalism. The bitter experience of One World united by post-Christian capitalism is the fuel which stokes the rationalizations of those who theorize it. So that the enemy is actually the post-Christian economic order, an order emanating from New York, London, and Berlin. Steve Sailer calls this the ``invade the world-invite the world\footnote{\url{http://isteve.blogspot.com/2005/02/invade-world-invite-world-personified.html}}'' strategy. The true and ancient enemy is not the alien among us, but rather the mind behind the mask which connives and schemes and manipulates that presence, which wishes bad fortune to come upon the lower cast. As Mencius Moldbug\footnote{\url{http://unqualified-reservations.blogspot.com/2013/06/civil-liberties-and-single-reactionary.html}} puts it, the State Department has an interest in keeping the American military busy extending its boundaries over seas, rather than tending to home affairs.

But who to blame? Should one assign responsibility, if it cannot be found? Should one, in fact, infer it?

I do not believe that this is a noble course of action for the truly resolute, for those who follow the path of Hercules. Does Hercules concern himself over-much with Eurystheus and his plots? Does he rage against Hera of the jealous eyes? Does he mutter indignities against the powers-that-be?

He pursues the mark relentlessly, yet almost casually. He has made the burden light. This making the burden light is what those who prophecy gloom miss. They do not feel equal to the challenge. I feel obligated to point out that the Dark Age ethos which is so often feted by the New Right was in fact likely lead by a transmigratory elite of nobility, as the intrigues by Celtic chieftains against Julius Caesar often made clear. It is far from clear to me that what was happening wasn't actually a family affair, in the sense that various initiates, chiefs, warriors and their on-hangers engaged in the leadership of the various bands or tribes of men living in the North of Europe. The genetic and clannish unity was an accidental by-product of a culture focused on high Northerness. Indeed, this appears to have been the normative pattern throughout Europe until our own day\footnote{\url{http://en.wikipedia.org/wiki/Abram\_Gannibal}}, \& the pursuit of such actually gave an accidental advantage to those of ``Northerness'', en masse. The subtle ruled, and benefited, the dense. Without that subtlety, it is just more ``fruitless populism'', I fear.

All the same, it is good (as Jeeves says) to ``know what tune the devil is playing''. It is certainly wise and prudent to assess and measure and identify one's spiritual adversaries, to maintain a sort of casual and detached regard of their latest fads and pursuits, and to know their names. Has anyone forgotten how Italy provoked\footnote{\url{http://www.worldcupblog.org/world-football/world-cup-moments-zinedine-zidane-headbutts-marco-materazzi.html}} France's star player into throwing away the World Cup? It's useful to know weak points in the make up of the adversary.

So far, as yet, the Right remains on the defensive, primarily because it is simply impossible to combat point-by-point every single inanity and injustice which is forwarded from the Left. As the crisis deepens, and as up and down become more indistinguishable to all but a few, more opportunity will arise for those with proper preparation to begin taking active roles in providing the umbrella of order in a world raining brimstone.

Hercules is unmoved, both by the presence of his own Moriarty (replete with numerous henchmen), and by the presence of even preternatural intervention and conspiracy. In fact, he stands their worst blows, and turns the injustice upon its head. It is not personal, for him, but rather, an occasion for spoiling the Egyptians. Is he not the son of Zeus? What is coming, however, from his triumph, will lead to the complete vindication of every possible injustice, and upon every scale and in every aspect, which was thrown against him.

The path of heaven contains the plan for such a perfect vindication for those who are worthy of it. Since our enemy will have nothing but a fight, then victory will accidentally be part of this. It is said, in fencing, that one should look, not at the sword, but at the eyes, of the opponent. So I have heard it said... 

\flrightit{Posted on 2013-11-02 by Logres}

\begin{center}* * *\end{center}

\begin{footnotesize}
\begin{sffamily}
\texttt{jc00001 on 2013-11-03 at 21:21 said:}

Brilliant post, answers many questions about the condition of the world today from a Traditionalist perspective.


\hfill

\end{sffamily}
\end{footnotesize}

